\section{User and Company}

\begin{frame}[fragile]{\texttt{env.user}}
	\begin{itemize}
		\item \textbf{\texttt{env.user}} is the current \texttt{res.users} record.
\begin{lstlisting}
user = self.env.user
print(user.name)
print(user.login)
print(user.has_group('base.group_system'))
\end{lstlisting}
		\item Useful for:
		\begin{itemize}
			\item checking groups / permissions in Python
			\item writing the current user into a field
			\item personalizing business logic
		\end{itemize}
	\end{itemize}
\end{frame}

\begin{frame}[fragile]{\texttt{env.company} and \texttt{env.companies}}
	\begin{itemize}
		\item \textbf{\texttt{env.company}}: the active company.
		\item \textbf{\texttt{env.companies}}: allowed companies in the current environment.
\begin{lstlisting}
company = self.env.company
print(company.name)
print(self.env.companies.mapped('name'))
\end{lstlisting}
		\item Multi-company rules often depend on these values.
		\item Never hard-code company IDs when \texttt{env.company} is enough.
	\end{itemize}
\end{frame}

\begin{frame}[fragile]{Practical User / Company Patterns}
\begin{lstlisting}
def create(self, vals):
	vals.setdefault('company_id', self.env.company.id)
	vals.setdefault('user_id', self.env.user.id)
	return super().create(vals)
\end{lstlisting}
	\begin{itemize}
		\item Set defaults from the environment.
		\item Keep business documents linked to the correct company/user.
	\end{itemize}
\end{frame}
